\chapter{Literature Review}\label{ch:lit}

\section{Classical Approaches}

\textbf{Lakshmi \emph{et al.} (2016)}~\cite{lakshmi2016fuzzy} use a
one-against-all SVM with contextual and grammar-rule features for four
classes (Name, Organization, Place, Date). A fuzzy module resolves cases
where the SVM assigns a word to both Name and Organization, using whether
the sentence subject governs an action verb. On a 500-sentence in-house test
set they report an F-measure of 83.82\% for the SVM alone and 88.13\% with
the fuzzy module. The method is interpretable and needs little data, but it
relies on hand-crafted rules and a parser, and its private test set makes
the results impossible to reproduce. Sreeja and Pillai~\cite{sp2020towards}
survey the challenges of Malayalam NER and argue that subword
representations are needed for its rich morphology.

\section{Neural and Transformer Approaches}

BiLSTM-CRF models~\cite{huang2015bidirectional, lample2016neural} replaced
hand-crafted features with learned ones and use a CRF~\cite{lafferty2001conditional}
to decode the best tag sequence. They were later outperformed by pre-trained
Transformers~\cite{vaswani2017attention}:
\begin{itemize}
  \item \textbf{mBERT}~\cite{devlin2019bert}: trained on Wikipedia in 104
  languages. Malayalam is only a small share of its data.
  \item \textbf{IndicBERTv2}~\cite{doddapaneni2023towards}: BERT-base trained
  on IndicCorp~v2 (23 Indic languages and English), with a vocabulary built
  for Indic scripts.
  \item \textbf{MuRIL}~\cite{khanuja2021muril}: trained on 16 Indian
  languages and their transliterations.
  \item \textbf{Naamapadam / IndicNER}~\cite{mhaske2023naamapadam}: the
  largest public Indic NER corpus. Training labels are projected from
  English, and the test sets are human-annotated.
\end{itemize}
Gopalakrishnan \emph{et al.}~\cite{gopalakrishnan2025efficient} fine-tune
m-BERT for Malayalam and report 96.45\% F1. Their dataset and evaluation
protocol are not released.

\section{Base Paper}

Bharathi Mohan \emph{et al.}~\cite{mohan2023finetuned} (IEEE EASCT 2023)
fine-tune multilingual BERT with a softmax token-classification layer on
WikiAnn~\cite{pan2017cross} for Tamil, Malayalam, Marathi and Bengali.
Table~\ref{tab:base} gives their Malayalam results.

\begin{table}[H]
\centering
\caption{Malayalam results reported in the base paper~\cite{mohan2023finetuned} (WikiAnn).}
\label{tab:base}
\begin{tabular}{lccc}
\toprule
Entity & Precision (\%) & Recall (\%) & F1 (\%) \\
\midrule
LOC     & 79.8 & 91.26 & 85.50 \\
ORG     & 81.3 & 76.5  & 78.8  \\
PER     & 86.2 & 84.4  & 85.3  \\
Overall & 82.1 & 82.1  & 82.2  \\
\bottomrule
\end{tabular}
\end{table}

\textbf{Strengths.} A single multilingual encoder gives competitive NER in
four Indian languages without feature engineering.

\textbf{Gaps.}
\begin{enumerate}
  \item Each token is classified independently, so tag dependencies are
  ignored and invalid sequences (e.g.\ \texttt{O}$\rightarrow$\texttt{I-ORG})
  are possible.
  \item The paper does not describe how word labels are aligned with subword
  tokens, a key choice for agglutinative languages.
  \item It does not say whether scores are token-level or span-level, and
  some figures are inconsistent. For example, the Marathi ORG F1 (88.8\%)
  is higher than both its precision (84\%) and its recall (85.6\%).
  \item WikiAnn is annotated automatically and consists mostly of short
  title fragments rather than running text.
\end{enumerate}

This project addresses these gaps. It uses an Indic-specific encoder, makes
the subword alignment explicit, adds a CRF layer, and evaluates with strict
span-level metrics on the human-annotated Naamapadam test set.
